\documentclass[12pt,a4paper,violet,openany]{bbe}

\setlength{\parindent}{0pt}
\setlength{\parskip}{0.5em}

\usepackage{fontspec}
\IfFontExistsTF{Roboto Bold}{
    \setmainfont{Roboto}[
        BoldFont = {Roboto Bold},
        ItalicFont = {Roboto Italic},
        BoldItalicFont = {Roboto Bold Italic}
    ]
    \setsansfont{Roboto}[
        BoldFont = {Roboto Bold},
        ItalicFont = {Roboto Italic},
        BoldItalicFont = {Roboto Bold Italic}
    ]
}{
    \setmainfont{Arial}
    \setsansfont{Arial}
}
\renewcommand{\familydefault}{\sfdefault}

\usepackage[spanish]{babel}
\addto\captionsspanish{\renewcommand{\tablename}{Tabla}}
\decimalpoint
\usepackage{hyperref}
\usepackage{booktabs}
\usepackage{array}
\usepackage{enumitem}

% Forzar numeracion plana de secciones (1, 2, 3...) en clase book
\renewcommand{\thesection}{\arabic{section}}
\renewcommand{\thesubsection}{\thesection.\arabic{subsection}}

\title{Taller practico SIPREH}
\authors{Grupo de Investigaci\'on en Ingenier\'ia de los Recursos H\'idricos --- GIREH\\Facultad de Ingenier\'ia, Universidad Nacional de Colombia}
\date{\today}

\begin{document}
\maketitle

\section{Objetivo del taller}
Aplicar las principales funcionalidades de SIPREH para consultar, explorar, comparar e interpretar informaci\'on relacionada con sequ\'ias meteorol\'ogicas, sequ\'ias hidrol\'ogicas y predicciones, mediante el an\'alisis de series temporales y mapas espaciales.

\section{An\'alisis hist\'orico}

\subsection{Sequ\'ias meteorol\'ogicas}
\begin{enumerate}[leftmargin=*,label=\alph*)]
    \item Genere una serie temporal de la variable clim\'atica de precipitaci\'on con resoluci\'on mensual para la celda ubicada en las coordenadas (-74.425, 3.675), localizada en el extremo sur, con una resoluci\'on espacial de 0,05\textdegree{}. Considere como periodo de an\'alisis el comprendido entre el 01/01/2007 y el 31/01/2008. A partir de la serie temporal obtenida, identifique los valores m\'aximo y m\'inimo de precipitaci\'on.
    \item Genere un mapa espacial del \i ndice SPI, con nivel de agregaci\'on de 3 meses y unidad de espacial de cuencas, utilizando como fuente de datos CHIRPS (0,05\textdegree{}), para el periodo comprendido entre el 01/01/2016 y el 29/02/2016. A partir de los resultados obtenidos, determine el grado de severidad de la sequ\'ia para la cuenca del Chuza.
    \item Cambie el tipo de visualizaci\'on de mapa espacial a serie temporal y modifique el nivel de agregaci\'on de 3 meses a 1 mes. En el periodo de analisis indique desde 01/01/2016 a 31/12/2016 para la cuenca Chisaca. Identifique el valor m\'inimo, valor m\'aximo, valor promedio del indice SPI.
    \item Genere un mapa del \i ndice RAI, con nivel de agregaci\'on de 6 meses y unidad espacial de celda, utilizando como fuente de datos CHIRPS (0,05\textdegree{}) para la fecha 15/02/2016. Exporte la imagen generada. Modifique la fecha de consulta al 15/08/201. Exporte la imagen generada y compare los resultados.
\end{enumerate}

\subsection{Sequ\'ias hidrol\'ogicas}
\begin{enumerate}[leftmargin=*,label=\alph*)]
    \item Genere una serie temporal del \i ndice SDI, con un nivel de agregaci\'on de 1 mes, para el periodo comprendido entre el 01/01/2019 y el 31/12/2022, correspondiente a la estaci\'on R\'io Tunjuelito \textendash{} Puente Bosa. A partir de la serie obtenida, determine el valor m\'aximo del \i ndice.
    \item Cambie el tipo de visualizaci\'on a mapa espacial y la fecha  01/05/2016. Determine el grado de severidad de la sequ\'ia hidrol\'ogica para la estaci\'on R\'io Guajaro Centro.
\end{enumerate}

\section{Predicci\'on actual}
\begin{enumerate}[leftmargin=*,label=\alph*)]
    \item Genere un mapa espacial, a unidad espacial de cuenca, del \i ndice SPI, con nivel de agregaci\'on de 12 meses y horizonte de predicci\'on de 2 meses. A partir de los resultados obtenidos, determine el grado de severidad del \i ndice SPI para la cuenca Chisac\'a.
    \item Genere una serie temporal del \i ndice SPI, con unidad espacial de perimetro urbano y nivel de agregaci\'on de 3 meses. Seleccione el pol\'igono en el mapa y determine los valores m\'inimo, m\'aximo y promedio del \i ndice. Exporte los resultados en formatos JSON y PNG.
    \item Genere un mapa, a unidad de celda, del \i ndice SPI, con nivel de agregaci\'on de 3 meses y horizonte de predicci\'on de 3 meses. Determine el grado de severidad de la sequ\'ia para la celda ubicada en las coordenadas -74.075, 5.275) (Esquina superior derecha).
    \item Genere el an\'alisis de la anomal\'ia mediante la opci\'on \textquotedblleft Ver anomal\'ia 2D\textquotedblright. Para la celda de estudio, determine la anomal\'ia de precipitaci\'on y el valor promedio de precipitaci\'on correspondiente. Exporte los resultados en formatos JSON y PNG.
    \item Genere el an\'alisis de anomal\'ias con nivel de agregaci\'on de 1 mes y  horizonte de predicci\'on de 5 meses. A partir de los resultados obtenidos, identifique la anomal\'ia de precipitaci\'on para la celda de estudio e interprete su comportamiento.
\end{enumerate}

\section{Hist\'orico de predicciones}
\begin{enumerate}[leftmargin=*,label=\alph*)]
    \item Genere una visualizaci\'on de serie 1D, a unidad de cuenca, del \i ndice SPI, con nivel de agregaci\'on de 3 meses, para la cuenca San Rafael. Determine los valores m\'inimo y m\'aximo del \i ndice.
    \item Genere una visualizaci\'on 2D, a unidad espacial de celda, del \i ndice SPI, con nivel de agregaci\'on de 12 meses y horizonte de predicci\'on de 7 meses. Exporte la imagen generada.
\end{enumerate}

\end{document}
